\documentclass[11pt]{article}
\usepackage[a4paper,margin=1in]{geometry}
\usepackage{amsmath,amssymb,booktabs,array}
\usepackage{hyperref}
\usepackage{tikz}
\usetikzlibrary{arrows.meta,positioning}

\title{Methodology: EV Motor Remanufacturing Inventory Simulation}
\author{}
\date{}

\begin{document}
\maketitle
\tableofcontents
\bigskip

% ─────────────────────────────────────────────────────────────────────────────
\section{Overview}

This document describes the discrete-event simulation (DES) methodology used to
study inventory management for a \textbf{closed-loop remanufacturing system} for
EV motors.  The system has two replenishment channels: an internal
\emph{remanufacturing} channel that recovers returned cores, and an external
\emph{supplier} channel that covers any remaining deficit.  Because both channels
are stochastic, an analytical model such as EOQ cannot capture the interaction
between them, making simulation the appropriate tool.

The simulator is built with \textbf{SimPy}, a Python process-based DES framework.
A parametric $(s,S)$ replenishment policy is tested over a 104-week horizon and
evaluated via service level, cost, and inventory trajectory.

% ─────────────────────────────────────────────────────────────────────────────
\section{Notation}

\begin{tabular}{lll}
\toprule
Symbol & Unit & Description \\
\midrule
$t$             & week       & simulation time \\
$I(t)$          & units      & stock on hand at time $t$ \\
$IP(t)$         & units      & pipeline-aware inventory position at time $t$ \\
$W(t)$          & units      & remanufacturing work-in-process (cores) at time $t$ \\
$O(t)$          & units      & supplier units currently on order at time $t$ \\
$D$             & units/week & constant customer demand per week \\
$p_r$           & ---        & core return rate (fraction of fulfilled demand) \\
$p_y$           & ---        & remanufacturing yield (fraction of returned cores) \\
$\tau_r$        & weeks      & fixed remanufacturing lead time \\
$s$             & units      & reorder point \\
$S$             & units      & order-up-to level \\
$\mu_L,\sigma_L$& weeks      & mean and std dev of supplier lead time \\
$c_p$           & \$/unit    & purchase cost per supplier unit \\
$c_o$           & \$/order   & fixed ordering cost per supplier order \\
$c_h$           & \$/unit$\cdot$week & holding cost per unit per week \\
\bottomrule
\end{tabular}

% ─────────────────────────────────────────────────────────────────────────────
\section{System Description}

\subsection{Closed-Loop Supply Chain}

Each week $D$ units of customer demand arrive.  Fulfilled units generate core
returns with probability $p_r$ each, which enter a remanufacturing pipeline.
After a lead time of $\tau_r$ weeks, each returned core is successfully
remanufactured with independent probability $p_y$.  Remanufactured units are
added back to stock.  An $(s,S)$ policy triggers external supplier orders
whenever the pipeline-aware inventory position falls at or below $s$.

\subsection{Material Flow}

At the end of each week $t$, let $f_t$ denote the number of units fulfilled and
$b_t$ the corresponding cores returned.  The quantities satisfy:
\begin{align}
  f_t &= \min(D,\; I(t^-)), \\
  b_t &\sim \mathrm{Binomial}(f_t,\; p_r).
\end{align}
After the reman lead time $\tau_r$, the good units produced from batch $b_t$ are:
\begin{equation}
  g_t \sim \mathrm{Binomial}(b_t,\; p_y).
\end{equation}
Supplier orders placed at time $t_0$ arrive at time $t_0 + L$, where:
\begin{equation}
  L \sim \max\!\bigl(0.1,\; \mathcal{N}(\mu_L, \sigma_L^2)\bigr).
\end{equation}
The lower bound of $0.1$ weeks prevents degenerate zero lead times.

% ─────────────────────────────────────────────────────────────────────────────
\section{Inventory Position}

The \textbf{stock on hand} $I(t)$ counts only units physically available.
The \textbf{inventory position} augments $I(t)$ with all pipeline commitments:
\begin{equation}
  \boxed{IP(t) \;=\; I(t) \;+\; O(t) \;+\; W(t)\,p_y}
  \label{eq:ip}
\end{equation}
where $W(t)\,p_y$ is the expected number of good units still in the
remanufacturing pipeline.  Using the expectation (rather than the realized count)
prevents the ordering policy from reacting to noise in the reman pipeline while
still accounting for the work already in progress.

% ─────────────────────────────────────────────────────────────────────────────
\section{Replenishment Policy}

\subsection{(s, S) Order-Up-To Policy}

The system uses a \textbf{periodic-review $(s,S)$ policy} with review period
$T_r = 1$ week.  At each review instant:
\begin{enumerate}
  \item Compute $IP(t)$ via Equation~\eqref{eq:ip}.
  \item If $IP(t) \le s$, place a supplier order of size:
        \begin{equation}
          Q_t \;=\; \max\!\bigl(0,\;\lfloor S - IP(t) \rfloor\bigr).
          \label{eq:order_qty}
        \end{equation}
  \item Otherwise, take no action.
\end{enumerate}
The order quantity $Q_t$ is \emph{variable}: it restores the pipeline position
exactly to $S$, unlike a fixed-$Q$ policy.  The use of $IP(t)$ in
Equation~\eqref{eq:order_qty} means the policy automatically adjusts for
outstanding supplier deliveries and pending reman output, avoiding double-ordering.

\subsection{Baseline Parameter Values}

\begin{center}
\begin{tabular}{lrl}
\toprule
Parameter & Value & Description \\
\midrule
$s$              & 35 units & reorder point \\
$S$              & 60 units & order-up-to level \\
$D$              & 25 units/week & weekly demand \\
$p_r$            & 0.87 & core return rate \\
$p_y$            & 0.90 & remanufacturing yield \\
$\tau_r$         & 2.0 weeks & reman lead time \\
$\mu_L$          & 2.0 weeks & mean supplier lead time \\
$\sigma_L$       & 0.5 weeks & std dev of supplier lead time \\
$c_p$            & \$400/unit & purchase cost \\
$c_o$            & \$250/order & fixed ordering cost \\
$c_h$            & \$4/unit$\cdot$week & holding cost \\
Horizon          & 104 weeks & simulation duration \\
\bottomrule
\end{tabular}
\end{center}

\subsection{Effective Reman Contribution}

Under the baseline parameters, the average net remanufacturing contribution per
week is:
\begin{equation}
  \bar{g} \;=\; D \cdot p_r \cdot p_y
          \;=\; 25 \times 0.87 \times 0.90
          \;\approx\; 19.6 \text{ units/week}.
\end{equation}
The remaining average external supplier need is therefore:
\begin{equation}
  D - \bar{g} \;=\; 25 - 19.6 \;\approx\; 5.4 \text{ units/week}.
\end{equation}
The closed-loop channel thus covers approximately $78\%$ of weekly demand on
average, making the supplier an incremental top-up rather than the primary source.

% ─────────────────────────────────────────────────────────────────────────────
\section{Cost Model}

Total cost over the simulation horizon $T$ (in weeks) is decomposed into three
components:
\begin{equation}
  \boxed{TC \;=\; c_p \sum_i q_i^{\mathrm{sup}}
              \;+\; c_o \cdot N_{\mathrm{ord}}
              \;+\; c_h \cdot \bar{I} \cdot T}
  \label{eq:tc}
\end{equation}
where:
\begin{itemize}
  \item $q_i^{\mathrm{sup}}$ is the quantity delivered by supplier order $i$,
  \item $N_{\mathrm{ord}}$ is the total number of supplier orders placed,
  \item $\bar{I} = \frac{1}{|\mathcal{S}|}\sum_{t \in \mathcal{S}} I(t)$ is the
        time-average stock on hand, computed from all state snapshots
        $\mathcal{S}$.
\end{itemize}
Remanufacturing costs are not modelled explicitly; they are treated as sunk or
captured in a separate cost centre.

% ─────────────────────────────────────────────────────────────────────────────
\section{Key Performance Indicators}

\begin{center}
\begin{tabular}{lll}
\toprule
KPI & Formula & Notes \\
\midrule
Service level & $\mathrm{SL} = \dfrac{\sum_t f_t}{\sum_t D}$ & fill-rate metric \\[6pt]
Total stockout units & $\sum_t (D - f_t)^+$ & $(x)^+ = \max(0,x)$ \\[4pt]
Avg stock on hand & $\bar{I}$ & from snapshots \\[4pt]
Avg inventory position & $\overline{IP}$ & from snapshots \\[4pt]
Avg expected reman WIP & $\overline{W \cdot p_y}$ & from snapshots \\[4pt]
Total cost & $TC$ & Equation~\eqref{eq:tc} \\
\bottomrule
\end{tabular}
\end{center}

% ─────────────────────────────────────────────────────────────────────────────
\section{Simulation Engine}

\subsection{Discrete-Event Simulation}

The SimPy engine advances simulation time to the next scheduled event.  Four
concurrent coroutine processes run inside a single shared state object:

\begin{description}
  \item[\texttt{weekly\_demand}] Fires every 1 week.  Computes $f_t$, records
    any stockout $b_t^{\mathrm{short}} = D - f_t$, draws
    $b_t \sim \mathrm{Binomial}(f_t, p_r)$, and spawns a
    \texttt{reman\_process} coroutine for each non-zero batch.

  \item[\texttt{reman\_process}($b_t$)] Suspends for $\tau_r$ weeks, then draws
    $g_t \sim \mathrm{Binomial}(b_t, p_y)$ and adds $g_t$ to $I(t)$.

  \item[\texttt{reorder\_monitor}] Fires every $T_r$ weeks.  Evaluates
    $IP(t)$ via Equation~\eqref{eq:ip}; if $IP(t) \le s$, places a supplier
    order of size $Q_t$ (Equation~\eqref{eq:order_qty}) and spawns a
    \texttt{supplier\_order} coroutine.

  \item[\texttt{supplier\_order}($Q_t$)] Draws $L \sim \max(0.1,
    \mathcal{N}(\mu_L, \sigma_L^2))$, suspends for $L$ weeks, then adds $Q_t$
    to $I(t)$.
\end{description}

\subsection{State Snapshots}

After every state-changing event, the simulator appends the current values of
$I(t)$, $IP(t)$, and $W(t)\,p_y$ to time-series logs.  These logs are used
for computing $\bar{I}$ and $\overline{IP}$ and for visualisation.

\subsection{Random Number Generation}

All stochastic draws use NumPy's \texttt{default\_rng} generator seeded for
reproducibility.  Three independent distributions are drawn per event type:

\begin{center}
\begin{tabular}{llll}
\toprule
Variable & Distribution & Parameters & Event \\
\midrule
$b_t$ & $\mathrm{Binomial}$ & $n = f_t$, $p = p_r$ & after demand \\
$g_t$ & $\mathrm{Binomial}$ & $n = b_t$, $p = p_y$ & after reman LT \\
$L$   & $\mathcal{N}$ (clipped) & $\mu_L$, $\sigma_L$ & per supplier order \\
\bottomrule
\end{tabular}
\end{center}

% ─────────────────────────────────────────────────────────────────────────────
\section{Single-Run Results (Baseline)}

With seed 42 and the baseline parameters over 104 weeks:

\begin{center}
\begin{tabular}{lr}
\toprule
KPI & Value \\
\midrule
Total demand              & 2{,}575 units \\
Filled demand             & 2{,}230 units \\
Total stockouts           & 345 units (60 events) \\
Service level ($\mathrm{SL}$) & 86.6\% \\
Reman units restocked     & 1{,}728 \\
Supplier units delivered  & 442 \\
Supplier orders placed    & 17 \\
Avg order quantity        & 27.5 units \\
Avg stock on hand ($\bar{I}$) & 11.2 units \\
Avg inventory position ($\overline{IP}$) & 41.8 units \\
Avg expected reman WIP    & 23.4 units \\
\midrule
Purchase cost ($c_p \sum q_i^{\mathrm{sup}}$) & \$176{,}800 \\
Ordering cost ($c_o N_{\mathrm{ord}}$)        & \$4{,}250 \\
Holding cost ($c_h \bar{I} T$)                & \$4{,}648 \\
\textbf{Total cost} ($TC$)                    & \$\mathbf{185{,}698} \\
\bottomrule
\end{tabular}
\end{center}

The service level of $86.6\%$ indicates that the reorder point $s = 35$ is
insufficient: during the 2-week reman and supplier lead times, demand can deplete
stock before replenishment arrives.  Raising $s$ trades higher holding cost for
fewer stockouts.

% ─────────────────────────────────────────────────────────────────────────────
\section{Monte Carlo Extension}

A single run with one seed is subject to sampling noise.  The natural extension
is to run $N$ independent replications with different seeds for a given policy
$(s, S)$ and aggregate KPIs into empirical distributions.  Let
$\mathrm{SL}^{(k)}$ and $TC^{(k)}$ denote the service level and total cost from
replication $k$.  The Monte Carlo estimates are:
\begin{align}
  \hat{\mu}_{\mathrm{SL}} &= \frac{1}{N}\sum_{k=1}^{N} \mathrm{SL}^{(k)}, &
  \hat{\sigma}_{\mathrm{SL}} &= \sqrt{\frac{1}{N-1}\sum_{k=1}^{N}
    \bigl(\mathrm{SL}^{(k)} - \hat{\mu}_{\mathrm{SL}}\bigr)^2}.
\end{align}
A $95\%$ confidence interval for the true mean service level is approximately:
\begin{equation}
  \hat{\mu}_{\mathrm{SL}} \;\pm\; 1.96\,\frac{\hat{\sigma}_{\mathrm{SL}}}{\sqrt{N}}.
\end{equation}
Sweeping $s$ over a grid $\{s_1, s_2, \ldots\}$ while holding $S$ fixed, and
plotting $(\hat{\mu}_{TC},\; \hat{\mu}_{\mathrm{SL}})$ for each value, traces
the \textbf{cost--service Pareto frontier} of the $(s,S)$ policy family.

% ─────────────────────────────────────────────────────────────────────────────
\section{Limitations and Future Work}

\begin{itemize}
  \item \textbf{Demand model}: demand is currently constant at $D$ units/week.
    Replacing it with a stochastic process (e.g.\ Poisson or negative binomial)
    would increase realism and widen the stockout distribution.
  \item \textbf{Reman lead time}: $\tau_r$ is deterministic.  A random
    $\tau_r \sim \mathcal{N}(\mu_r, \sigma_r^2)$ would capture shop-floor
    variability.
  \item \textbf{Reman cost}: adding a cost per core ($c_r$ per unit processed)
    would allow comparison of reman-heavy vs.\ supplier-heavy strategies on a
    total-cost basis.
  \item \textbf{Lost sales vs.\ backlog}: the current model uses lost sales
    (unfulfilled demand is not backordered).  A backlog model would track
    cumulative shortfall and may be more appropriate if customers wait.
  \item \textbf{Capacity constraints}: the remanufacturing shop is assumed to
    have unlimited capacity.  A bounded throughput constraint would add queuing
    dynamics.
\end{itemize}

\end{document}
